\documentclass[11pt]{article}
\usepackage[T1]{fontenc}
\usepackage[margin=1in]{geometry}
\usepackage{amsmath,amssymb,amsthm,booktabs,graphicx,microtype}
\usepackage[hidelinks]{hyperref}
\newtheorem{theorem}{Theorem}[section]
\newtheorem{proposition}[theorem]{Proposition}
\newtheorem{corollary}[theorem]{Corollary}
\theoremstyle{definition}
\newtheorem{definition}[theorem]{Definition}
\newtheorem{remark}[theorem]{Remark}
\newcommand{\Aut}{\operatorname{Aut}}
\newcommand{\Out}{\operatorname{Out}}
\newcommand{\E}{\mathbb E}
\newcommand{\Q}{\mathbb Q}
\newcommand{\Var}{\operatorname{Var}}
\numberwithin{equation}{section}
\title{Cubic Refinements and Euler Characteristics of Outer Space}
\author{Alexander Kolpakov\\[0.5ex]
\small Center for Science, Technology, Engineering, and Mathematics\\
\small University of Austin, Austin, Texas, United States\\
\small \texttt{akolpakov@uaustin.org}\\[0.75ex]
\small The Wolfram Institute for Computational Foundations of Science\\
\small \texttt{akolpakov@wolframinstitute.org}}
\date{Research draft, 20 October 2026}
\begin{document}
\maketitle
\begin{abstract}
The stationary law of the local Whitehead chain on cubic multigraphs is
proportional to the reciprocal of the full incidence automorphism order.
We use cubic refinements to transfer finite sums over all stable
rank-$r$ graphs to expectations under this law at $2r-2$ vertices.
The refinement multiplicity at a vertex of valence $d$ is $(2d-5)!!$.
Applying this identity to the established forest formulas gives distinct
observables for the ordinary and virtual Euler characteristics of
$\Out(F_r)$. We explain the orientation and stabilizer data required by
the forested chain complex, and give exact low-rank checks, variances,
and a two-term representative of the known first Morita class.
The transfer does not give an efficient sampling algorithm: the present
observables sum forests, and their stationary variances omit
autocorrelation. The topology and Euler formulas used here are known;
the purpose is to relate them precisely to the cubic walk.
\end{abstract}

\section{Introduction}
A Whitehead move on a cubic graph collapses an edge and expands the
resulting four-valent vertex in a different way. This is also a basic
local operation in Outer space, but the finite graph of cubic
isomorphism classes records only part of its structure. It forgets
markings, edge lengths, forests, and orientation characters.
These omissions matter both for homology and for Euler characteristics.

Let $\mathcal G_N$ be the connected cubic multigraphs on $N$ vertices,
up to isomorphism, with loops and parallel edges allowed. The local
chain of \cite{Kolpakov} chooses an edge and one of its two alternative
resolutions, counting moves with multiplicity. A choice of a loop
leaves the graph unchanged. Its stationary distribution is
\begin{equation}\label{eq:stationary}
 \pi_N(G)=\frac{1}{C_N|\Aut(G)|},\qquad
 C_N=\sum_{[G]\in\mathcal G_N}\frac1{|\Aut(G)|}.
\end{equation}
Here and below $\Aut(G)$ acts on vertices and half-edges. Loop flips
and permutations of parallel edges belong to this group.

For a fixed rank $r\ge2$, every connected stable graph is a contraction
of a cubic graph on $N=2r-2$ vertices. Counting those contractions
without their multiplicities gives an incorrect sum. Our main
identity restores the multiplicities explicitly.

\begin{theorem}[Refinement transfer]\label{thm:transfer}
Let $\mathcal S_r$ be the connected rank-$r$ multigraphs with all
valences at least three. For $H\in\mathcal S_r$ put
\[
 R(H)=\prod_{v\in V(H)}(2\deg(v)-5)!!.
\]
For every isomorphism-invariant function $u:\mathcal S_r\to\Q$,
\begin{equation}\label{eq:transfer}
 \sum_{[H]\in\mathcal S_r}\frac{u(H)}{|\Aut(H)|}
 =C_N\E_{\pi_N}\!
   \left[\sum_{\substack{F\subseteq E(G)\\F\ {\rm forest}}}
                     \frac{u(G/F)}{R(G/F)}\right].
\end{equation}
Each subset of parallel edges is counted separately in the inner sum.
Loops cannot belong to a forest.
\end{theorem}

The proof is an orbit--stabilizer count, an instance of the established
groupoid-cardinality principle \cite{Groupoidification}.
The identity is useful because
the right side uses a single fixed cubic state space while retaining
all stable strata. We make no priority claim for the underlying
refinement count. The following sections apply it to established
topological formulas and make the associated finite checks reproducible.

Culler--Vogtmann introduced Outer space in \cite{CV86}.
Conant--Vogtmann describe the relevant forested graph complexes in
\cite{ConantVogtmann}. The virtual Euler generating function originates
with Smillie--Vogtmann \cite{SmillieVogtmann}; Borinsky--Vogtmann develop
its asymptotics and formal-integral description in
\cite{BV2020,BVQFT}. The ordinary Euler characteristic and its generating
function are treated separately in \cite{BV2023}. None of these formulas,
or the low-rank homology groups reported below, is proposed as a new
topological invariant.

\section{Cubic chambers and forested cells}

\subsection{What a move records}
A rank-$r$ cubic graph has $2r-2$ vertices and $3r-3$ edges.
Positive edge lengths with total length one give a metric simplex
of dimension $3r-4$. A marking identifies its fundamental group with
$F_r$. Forgetting both the marking and the lengths gives a state of
the finite walk.

The spine of Outer space has dimension $2r-3$. If $e$ is a non-loop
edge and a Whitehead move changes $G$ to $G'$, it records a span
\[
 G\longrightarrow G/e\longleftarrow G'.
\]
This span is a path of two edges in the spine. Higher-dimensional
cells require several compatible forest collapses. Chamber adjacency
alone therefore does not recover the spine's chain complex.

\subsection{Orientations and stabilizers}
Let $C_k(r)$ be the rational vector space generated by pairs $(G,F)$,
where $G\in\mathcal S_r$ and $F$ is a forest of $k$ ordered edges.
An odd permutation of the forest order changes the sign, and graph
isomorphisms identify generators. In particular, a forest whose
stabilizer contains an odd permutation represents zero.

For the order $e_0,\ldots,e_{k-1}$ define
\begin{align}\label{eq:differential}
 \partial_c(G,F)&=\sum_i(-1)^i(G/e_i,F/e_i),\\
 \partial_u(G,F)&=\sum_i(-1)^i(G,F\setminus e_i),\qquad
 \partial=\partial_c-\partial_u.\nonumber
\end{align}
The remaining edge order is inherited. Two successive contractions
cancel in opposite orders, as do two unmarkings. A contraction and an
unmarking anticommute. Thus
\[
 \partial_c^2=\partial_u^2
 =\partial_c\partial_u+\partial_u\partial_c=\partial^2=0.
\]
With finite stabilizers treated in characteristic zero, this is the
forested cubical model computing $H_*(\Out(F_r);\Q)$
\cite[Section 3.2]{ConantVogtmann}. The presentation by cubic graphs
modulo IHX arises from a further filtration; it is not the same
generator list.

The full incidence group is essential. If $l_v$ is the number of
loops at $v$ and $m_{uv}$ the number of edges between distinct vertices,
then
\begin{equation}\label{eq:aut}
 |\Aut(G)|=|\Aut_{\rm vertices}(G)|
               \prod_v2^{l_v}l_v!\prod_{u<v}m_{uv}!.
\end{equation}
If a forest uses one edge of a parallel class, the corresponding
factor in its stabilizer is divided by $m_{uv}$, after restricting
the vertex group to preserve the marked pairs. A forest uses at most
one edge from each parallel class. One may quotient those choices when
building a chain basis, but must restore their multiplicities in
stabilizer-weighted sums.

At rank two, the rose with two loops has automorphism order eight.
The marked theta and dumbbell contract to it with coefficient one
in the oriented orbit basis. In the stabilizer pairing
$\langle [G,F],[G,F]\rangle=|\Aut(G,F)|$, the adjoint contraction sends
the rose to $2\,\theta+\mathrm{dumbbell}$. An ordinary matrix transpose
would give coefficients one and one. Move multiplicities and cellular
boundary coefficients therefore require different pairings.

\section{The refinement identity}

\begin{proposition}\label{prop:trees}
A trivalent tree with $d$ distinctly labelled incident leaves has
$(2d-5)!!$ isomorphism classes.
\end{proposition}
\begin{proof}
There is one tree for three leaves. Given a tree with $d-1$ leaves,
subdivide any of its $2d-5$ edges and attach the last labelled leaf.
Removing that leaf and suppressing the resulting bivalent vertex is
the unique inverse operation. Induction gives the formula.
\end{proof}

\begin{proof}[Proof of Theorem~\ref{thm:transfer}]
Fix a stable graph $H$. Replace each vertex of valence $d$ by a
trivalent tree with its incident half-edges as distinct leaf labels.
Proposition~\ref{prop:trees} gives $R(H)$ choices. The inserted tree
edges form a forest $F$ in a cubic graph $G$, and $G/F\cong H$.
Conversely every such pair $(G,F)$ determines one of these refinements
after choosing an identification of its contraction with $H$.

The group $\Aut(H)$ acts on the refinement choices, and their orbits
are precisely the isomorphism classes of these pairs. The stabilizer
of a refinement is $\Aut(G,F)$: an automorphism of a trivalent tree
fixing all its incident labels is trivial. Orbit--stabilizer gives
\begin{equation}\label{eq:fibre}
 \sum_{\substack{[(G,F)]\\G/F\cong H}}\frac1{|\Aut(G,F)|}
              =\frac{R(H)}{|\Aut(H)|}.
\end{equation}
For each fixed cubic $G$, summing over its forest orbits with weight
$1/|\Aut(G,F)|$ equals summing over all its individual forests with
weight $1/|\Aut(G)|$. Multiply \eqref{eq:fibre} by $u(H)/R(H)$ and
sum over $H$. Equation~\eqref{eq:stationary} proves
\eqref{eq:transfer}. All sums are finite.
\end{proof}

The correction can be calculated directly on a cubic graph.
A component of $F$ with $s$ vertices contracts to a vertex of valence
$3s-2(s-1)=s+2$. Consequently
\begin{equation}\label{eq:components}
 R(G/F)=\prod_{K\in\pi_0(F)}(2|V(K)|-1)!!,
\end{equation}
including isolated vertices as forest components.
No graph automorphism calculation is needed for this denominator.

\begin{remark}
The theorem applies to any stable-graph observable, including one
which depends on a known automorphism order or orientation character.
It does not remove the cost of evaluating that observable.
In particular, its use in an expectation does not by itself give a
polynomial-time randomized algorithm.
\end{remark}

\section{The two Euler characteristics}

\subsection{The virtual invariant}
Write
\[
 f_H(-1)=\sum_{\substack{F\subseteq E(H)\\F\ {\rm forest}}}(-1)^{|F|}.
\]
The established forest formula is
\begin{equation}\label{eq:virtual}
 \chi_{\rm virt}(\Out(F_r))
       =\sum_{[H]\in\mathcal S_r}\frac{f_H(-1)}{|\Aut(H)|};
\end{equation}
see \cite[Section 4]{BVQFT}. The sign records cube dimension, while
the denominator records its finite stabilizer.
For a connected graph the Tutte subset formula gives
\[
 f_H(-1)=(-1)^{|V(H)|-1}T_H(0,1).
\]
If $H$ has a bridge, adding or removing that bridge pairs forests of
opposite signs, so $f_H(-1)=0$.

Theorem~\ref{thm:transfer} now gives the following complete expectation.
\begin{corollary}\label{cor:virtual}
For $N=2r-2$,
\[
 \chi_{\rm virt}(\Out(F_r))=C_N\E_{\pi_N}X(G),\qquad
 X(G)=\sum_{F\ {\rm forest\ in}\ G}\frac{f_{G/F}(-1)}{R(G/F)}.
\]
\end{corollary}
By contrast, $-C_N\E_{\pi_N}T_G(0,1)$ gives only the cubic contribution
to \eqref{eq:virtual}. Lower-vertex strata must be restored.

\subsection{The ordinary invariant}
Let $a_k(H)$ be the number of $\Aut(H)$-orbits of $k$-edge forests
whose stabilizer contains no odd permutation of the forest. Set
\begin{equation}\label{eq:epsilon}
 \epsilon(H)=\sum_k(-1)^k a_k(H).
\end{equation}
The oriented quotient chain complex gives
\[
 e(\Out(F_r))=\sum_{[H]\in\mathcal S_r}\epsilon(H).
\]
This is the ordinary rational Euler characteristic. It is generally
different from the virtual one.

For clarity, the character form of Burnside's formula reads
\begin{equation}\label{eq:burnside}
 \epsilon(H)=\frac1{|\Aut(H)|}
  \sum_{\alpha\in\Aut(H)}
   \sum_{\substack{F\ {\rm forest}\\\alpha F=F}}
                         (-1)^{c(\alpha|_F)},
\end{equation}
where $c$ counts cycles of the permutation on the forest edges.
Indeed the trace of $\alpha$ on the oriented forest basis is
$\operatorname{sgn}(\alpha|_F)$. Multiplying by $(-1)^{|F|}$ gives
$(-1)^{c(\alpha|_F)}$. Averaging projects to the invariant subspace.
This is the formula of \cite[Theorem 2.1]{BV2023}, grouped here by
the underlying graph.

The identity automorphism contribution to the sum over $H$ is
\eqref{eq:virtual}. The other automorphisms supply the difference
between the invariants. If parallel choices have first been
quotiented out, the identity of the remaining vertex group alone
does not recover that contribution.

Apply Theorem~\ref{thm:transfer} to
$u(H)=|\Aut(H)|\epsilon(H)$.
\begin{corollary}\label{cor:ordinary}
For $N=2r-2$,
\[
 e(\Out(F_r))=C_N\E_{\pi_N}Y(G),\qquad
 Y(G)=\sum_{F\ {\rm forest\ in}\ G}
       \frac{|\Aut(G/F)|\epsilon(G/F)}{R(G/F)}.
\]
\end{corollary}
The factors $|\Aut(G/F)|$ cannot be omitted: the target sum over
stable graph types is unweighted.

\subsection{Rank two}
There are three stable graph types: the rose with two loops, the theta,
and the dumbbell. Their automorphism orders are $8,12,8$.
The respective forest sums are $1,-2,0$, so
\[
 \chi_{\rm virt}(\Out(F_2))=\frac18-\frac2{12}=-\frac1{24}.
\]
The rose has $\epsilon=1$ and the other two types have $\epsilon=0$,
giving $e(\Out(F_2))=1$.

The cubic law has $C_2=5/24$ and probabilities $(2/5,3/5)$ on
the theta and dumbbell. The corrected observables are
\[
 (X(\theta),X(\mathrm{dumbbell}))=(-1,1/3),\qquad
 (Y(\theta),Y(\mathrm{dumbbell}))=(8,8/3).
\]
Their expectations give both invariants with the same cubic law.
This example detects missing parallel-edge choices, missing loop
flips and the confusion between the two Euler characteristics.

\section{Exact computations and their scope}

\subsection{Independent enumeration checks}
Every stable graph is reached by repeatedly contracting non-loop
edges in a cubic refinement. The programs enumerate those types,
deduplicate them, and compute the signed forest sum.
For each stable type they separately verify the fibre identity
\eqref{eq:fibre} using exact rational arithmetic.

An independent check uses the formal Gaussian generating integral
of \cite{BVQFT}. If $Z$ is a standard real Gaussian, put
\[
 T(x)=-x-x^2/2+(1+x)\log(1+x),\qquad
 \exp\mathcal Y(\hbar)=
       \E\exp\{-\hbar^{-1}T(\sqrt\hbar Z)\}.
\]
Then $[\hbar^{r-1}]\mathcal Y=\chi_{\rm virt}(\Out(F_r))$.
The checker computes the same coefficients from vertex-degree profiles
and Gaussian moments, without enumerating graphs or forests.

\begin{center}
\begin{tabular}{rrrrr}
\toprule
Rank & Stable types & Cubic types & $\chi_{\rm virt}$ & $e$\\
\midrule
2 & 3 & 2 & $-1/24$ & 1\\
3 & 15 & 5 & $-1/48$ & 1\\
4 & 111 & 17 & $-161/5760$ & 2\\
5 & 1076 & 71 & $-367/5760$ & 1\\
\bottomrule
\end{tabular}
\end{center}
These reproduce known values. The virtual fibre check extends through
rank six, with 388 cubic classes and 2,746,060 individual forests.
The ordinary values are also checked against the alternating dimensions
of the oriented chain complex.

\subsection{Sampling variance}
For independent samples from $\pi_N$, the relative variances are:
\begin{center}
\begin{tabular}{rrr}
\toprule
Rank & $\Var(X)/(\E X)^2$ & $\Var(Y)/(\E Y)^2$\\
\midrule
2 & 10.667 & 0.29630\\
3 & 17.052 & 0.50612\\
4 & 20.452 & 3.96064\\
5 & 21.284 & 74.40849\\
\bottomrule
\end{tabular}
\end{center}
The observables estimate different targets, so the table does not
compare two algorithms for the same invariant. The calculations also
omit Markov-chain autocorrelation and the cost of the forest sum.
The small-rank values prove no asymptotic variance bound.

\subsection{A concrete cycle certificate}
The full oriented chain dimensions through rank four are
\[
 (3,2),\qquad (15,22,9,1),\qquad (111,310,385,312,181,53).
\]
Exact rational ranks give $H_0=\Q$ in ranks two and three and
$H_0=H_4=\Q$ in rank four, with other groups zero.
The separate rank-five certificate also recovers $H_0=\Q$ and no
positive-degree homology. These are established low-rank results.

One compact representative is worth recording. On the triangular
prism with triangles $012$ and $345$ and matching edges $03,14,25$,
let the ordered forests be
\[
 A=(01,02,34,35),\qquad B=(01,12,34,35),\qquad z=A-2B.
\]
\begin{figure}[ht]
\centering
\includegraphics[width=.82\textwidth]{prism-cycle.pdf}
\caption{The two ordered forests in the prism representative.}
\end{figure}
For a triangle with ordered edges $a,b,c$, the forest chain
$t=a\wedge b-a\wedge c+b\wedge c$ has $\partial_u t=0$.
The product for the two triangles has nine terms. In the oriented
prism quotient, three terms give $3A$ and six give $-6B$.
The attached certificate verifies
\[
 \partial_c z=\partial_u z=0,\qquad
 \phi\partial_5=0,\qquad \phi(z)=4
\]
for an integer cochain $\phi$ with 18 nonzero terms. Thus $z$ is
not a boundary. The serialized graphs, edge orders and coefficients
allow an independent check without trusting a rank or nullspace search.
This is the known first Morita class: the prism is obtained by
expanding the vertices of the three-edge theta graph into circles,
as in \cite[Section 6.1]{Morita}. It is not a new homology class.

\section{What remains open}
The refinement identity transfers exact sums, but its usefulness for
large rank depends on the cost and variance of the resulting observable.
A useful next result would control one of those quantities, or replace
the forest sum by a cheaper estimator with a proved error bound.

For homology, restricting to three-edge-connected graphs substantially
reduces the finite complexes. Grego--V{\'i}tek--Willwacher
\cite[Appendix A]{GVW} distinguish the full differential from the
contracting differential. Their Conjecture 8 states that contracting
homology in the three-edge-connected complex is concentrated in excess
zero, and is checked there through rank seven. Our smaller computations
do not extend that range. A general contraction argument must control
separating pairs created by refinements; an argument for unrestricted
vertex expansions does not automatically preserve three-edge-connectivity.

The general nonvanishing of Morita classes is no longer an open
target: see Kupers--Miller--Patzt \cite{KMP}. The present two-term
certificate concerns a known low-rank class. These questions are
independent of the quartic period torsor and its series realizations.

\subsection*{Reproduction}
The accompanying repository contains the source, exact archives and
standalone checks in \path{research/outer-forest-complex}.
The commands \texttt{make outer-space-paper} and
\texttt{make outer-space-check} build the manuscript and run the
default finite checks. The Python and native graph canonicalizers
are compared independently. Rational ranks are computed exactly
through rank four. The rank-five archive states separately where
modular rank bounds and acyclic quotient arguments certify the
rational answer. Agreement modulo two primes alone is not used as
a rational-rank certificate.

\bibliographystyle{abbrvurl}
\bibliography{outer-space}
\end{document}
